\chapter{Conclusions and Future Work}
\label{chap:conclusions}

\section{Conclusions}
This thesis presented FAIRE, a precision-oriented framework for force-aware imitation with reactive execution in robotic polishing. Its first component is a VR-calibrated handheld teaching system that transfers synchronized position, orientation, wrench, and image demonstrations into the robot workspace. Its second component is a force-aware imitation policy whose GRU encoder represents recent six-axis wrench history and whose output contains both motion references and a desired force trajectory. Its third component is an execution stack in which admittance control realizes the learned force reference through compliant Cartesian correction and bounded residual policies adjust sliding velocity and tool-normal alignment.

The framework makes a deliberate separation of responsibilities. The imitation policy defines task-level motion and contact references. The admittance controller converts force error into compliant motion. Residual RL supplies precision corrections without replacing the task. This separation makes the learned force output physically interpretable and enables controlled ablations of observation, action, representation, controller, and residual policy.

The planned experiments evaluate force tracking, contact stability, normal alignment, polishing-quality improvement, spatial consistency, and task time on curved headlight surfaces. The source FAIRE manuscript does not yet provide numerical results. Therefore, no quantitative superiority claim is made in this version of the thesis.

\section{Limitations}
The principal limitation is evidentiary completeness. Hardware identifiers, calibration accuracy, synchronization details, network dimensions, learning hyperparameters, admittance parameters, residual RL algorithms, safety bounds, trial counts, and quantitative results remain TBD. The current manuscript defines the intended system and experimental interfaces but is not yet sufficient for independent replication.

The framework also assumes a specified target polishing region and does not solve autonomous defect detection. Its demonstrations and policies may not generalize to unseen geometries, materials, tools, abrasives, or contact regimes. A predicted desired force does not guarantee safe physical execution; performance depends on sensor quality, frame consistency, controller tuning, saturation, and independent safety logic. Surface-normal estimation error can degrade orientation correction, and the velocity and orientation residuals may interact.

\section{Future Work}
Immediate work should complete the implementation record and populate the preregistered ablations with repeated real-robot trials. The calibration transform should be validated quantitatively, multimodal delay should be measured, and the discrete admittance implementation should be tested under conservative force limits before learned residuals are enabled.

The evaluation should then expand across multiple curved surfaces, target regions, tools, abrasive conditions, and initial-contact offsets. Surface quality should be measured with a repeatable instrument or controlled vision protocol and related to the tracked force, velocity, and alignment histories. Reporting interventions, saturation, and worst-case force is as important as reporting average quality.

Longer-term extensions include uncertainty-aware force prediction, safety-filtered residual actions, online contact-frame estimation, and adaptation across workpiece geometry and tool wear. Autonomous task-region localization may be integrated as a separate upstream module, but it should remain distinguishable from FAIRE's core contribution: faithful, compliant reproduction and precision correction of demonstrated polishing contact behavior.
